%% GENERATED by tools/gen_latex.py -- edit content/ch02.py instead.
%% Build:  latexmk -pdf ch02.tex     (or: pdflatex ch02.tex, twice)
\documentclass[aspectratio=169,11pt,t]{beamer}
\usepackage{itbpro}

\coursetitle{Blok Bangunan Matematis Jaringan Saraf}
\coursesubtitle{Tensor, operasi tensor, dan penurunan berbasis gradien -- dijelaskan lewat kode yang bisa dijalankan, bukan notasi.}
\coursekicker{Chapter 2}
\coursesource{Chollet \& Watson, Deep Learning with Python 3e -- bab 2}
\courseduration{3 jam (2 sesi)}
\coursepresenter{Rahman Indra Kesuma, S.Kom., M.Cs.}{Asisten Pengajar}
\coursebrand{AI for Professional \textperiodcentered\ ITB \textperiodcentered\ Ch 2}

\title{Blok Bangunan Matematis Jaringan Saraf}
\author{Rahman Indra Kesuma, S.Kom., M.Cs.}
\date{}

\begin{document}
\covertitle

\framekicker{Learning outcomes}
\begin{frame}[shrink=25]{Tujuan Sesi}
\leadpar{Setelah sesi ini, peserta dapat:}
\begin{isteps}
\item Menjalankan contoh MNIST pertama dari ujung ke ujung dan menyebut peran \textbf{layer, loss, optimizer, dan metric} pada tiap barisnya.
\item Membaca \textbf{rank, shape, dan dtype} sebuah tensor, dan memetakan data nyata (vektor, deret waktu, citra, video) ke bentuk tensornya.
\item Menjelaskan \textbf{operasi elemen-demi-elemen, broadcasting, hasil kali tensor, dan reshape} beserta arti geometrisnya.
\item Menerangkan \textbf{turunan, gradien, SGD mini-batch, dan aturan rantai}, lalu menunjuk letaknya di dalam graf komputasi.
\item Menulis ulang MNIST \textbf{dari nol} -- Dense, Sequential, batch generator, dan lingkar pelatihan -- tanpa memakai \inlinecode{fit()}.
\end{isteps}

\vskip 6pt
\vskip 4pt
\reslink{Buku}{Bab 2 --- teks penuh}{https://deeplearningwithpython.io/chapters/chapter02\_mathematical-building-blocks}
\reslink{Notebook}{01\_mnist\_pertama.ipynb}{../../notebooks/ch02/01\_mnist\_pertama.ipynb}
\reslink{Notebook}{02\_tensor\_dan\_operasi.ipynb}{../../notebooks/ch02/02\_tensor\_dan\_operasi.ipynb}
\reslink{Notebook}{03\_gradien\_dan\_sgd.ipynb}{../../notebooks/ch02/03\_gradien\_dan\_sgd.ipynb}
\reslink{Notebook}{04\_mnist\_dari\_nol.ipynb}{../../notebooks/ch02/04\_mnist\_dari\_nol.ipynb}
\reslink{Repo}{Notebook resmi penulis}{https://github.com/fchollet/deep-learning-with-python-notebooks/blob/master/chapter02\_mathematical-building-blocks.ipynb}
\vskip 2pt
\end{frame}

\framekicker{Peta bab}
\begin{frame}[shrink=25]{Satu contoh, dibongkar sampai ke bawah}
\leadpar{Bab 2 bergerak dalam satu lingkaran penuh: jalankan MNIST dengan \inlinecode{fit()}, bongkar tiap potongnya, lalu \hilite{tulis ulang seluruhnya dari nol} dan buktikan hasilnya sama.}
\begin{minipage}[t]{0.241\textwidth}
\begin{icard}
\cardh{1 $\cdot$ Contoh pertama}
MNIST dalam 10 baris. Berjalan dulu, dimengerti belakangan.
\cardtag{bag. 2.1}
\end{icard}
\end{minipage}%
\hfill
\begin{minipage}[t]{0.241\textwidth}
\begin{icard}
\cardh{2 $\cdot$ Tensor}
Rank, shape, dtype, irisan, sumbu batch, dan data dunia nyata.
\cardtag{bag. 2.2}
\end{icard}
\end{minipage}%
\hfill
\begin{minipage}[t]{0.241\textwidth}
\begin{icard}
\cardh{3 $\cdot$ Operasi tensor}
Elemen-demi-elemen, broadcasting, matmul, reshape, dan geometrinya.
\cardtag{bag. 2.3}
\end{icard}
\end{minipage}%
\hfill
\begin{minipage}[t]{0.241\textwidth}
\begin{icard}
\cardh{4 $\cdot$ Mesinnya}
Turunan, gradien, SGD, aturan rantai, autodiff.
\cardtag{bag. 2.4-2.6}
\end{icard}
\end{minipage}%
\begin{iquote}
\quotetext{Kode yang bisa dijalankan adalah keterangan paling tepat dan paling tidak ambigu untuk sebuah operasi matematis.}\quotecite{Semangat bab 2 -- notasi diganti implementasi}
\end{iquote}
\end{frame}
\note{Sesi ini panjang. Pecah di antara bagian 2.3 dan 2.4; separuh pertama soal data, separuh kedua soal belajar.}

\coursesection{01}{Sekali lihat jaringan saraf}{MNIST -- 'hello world'-nya deep learning.}

\framekicker{Bagian 2.1}
\begin{frame}[shrink=25]{Muat data: 60.000 latih, 10.000 uji}
\icode{python}{listing 2.1 --- memuat MNIST}{listings/ch02-01.py}
\ioutput{listings/ch02-02.txt}
\begin{itable}{>{\raggedright\arraybackslash}p{0.2256\textwidth}>{\raggedright\arraybackslash}p{0.7144\textwidth}}
\thead{Istilah} & \thead{Artinya di sini} \\ \midrule
\textbf{Sample} & Satu titik data -- satu citra 28$\times$28. \\
\textbf{Class} & Satu kategori -- angka 0 sampai 9. \\
\textbf{Label} & Kelas yang melekat pada satu sample tertentu. \\
\end{itable}
\end{frame}
\note{Tunjukkan satu citra dengan matplotlib sebelum lanjut; peserta perlu melihat bahwa 'data' di sini benar-benar gambar.}

\framekicker{Bagian 2.1}
\begin{frame}[shrink=25]{Model, kompilasi, pelatihan -- sepuluh baris}
\icode{python}{listing 2.2-2.5 --- MNIST ujung ke ujung}{listings/ch02-03.py}
\begin{iband}[signal]
Sebuah \textbf{layer} adalah \textit{penyaring data}: ia menerima data dan mengeluarkan representasi yang lebih berguna. \inlinecode{softmax} di lapis akhir mengeluarkan \hilite{10 skor peluang yang berjumlah 1}.
\end{iband}
\end{frame}
\note{Perhatikan: tidak ada input\_shape. Keras menyimpulkan bentuk masukan sendiri pada pemanggilan pertama --- dibahas di bab 3.}

\framekicker{Bagian 2.1}
\begin{frame}[shrink=25]{Hasilnya -- dan celah pertama yang harus dicurigai}
\ioutput{listings/ch02-04.txt}
\begin{minipage}[t]{0.485\textwidth}
\begin{minipage}[t]{0.494\textwidth}
\istat{98,9\%}{akurasi pada data latih}
\end{minipage}%
\hfill
\begin{minipage}[t]{0.494\textwidth}
\istat{97,8\%}{akurasi pada data uji}
\end{minipage}%

\end{minipage}%
\hfill
\begin{minipage}[t]{0.485\textwidth}
\begin{iband}[amber]
Selisih \textasciitilde{}1,1 poin itu bukan derau. Itu \textbf{overfitting}: model bekerja lebih baik pada yang pernah dilihatnya. Bab 5 membahasnya sebagai persoalan pokok.
\end{iband}

\end{minipage}%
\icode{python}{listing 2.6 --- meramal}{listings/ch02-05.py}
\ioutput{listings/ch02-06.txt}
\end{frame}
\note{Angka pastinya berayun tiap kali dilatih ulang --- katakan itu di depan, supaya peserta tidak panik kalau notebooknya memberi 97,6\%.}

\coursesection{02}{Representasi data: tensor}{Wadah untuk data. Tiga sifat, dan satu sumbu istimewa.}

\framekicker{Bagian 2.2}
\begin{frame}[shrink=25]{Rank, shape, dtype}
\begin{center}
\adjustbox{max width=\linewidth,max totalheight=0.52\textheight}{%

\begin{tikzpicture}[font=\sffamily\tiny,
  bx/.style={draw=signal!60, fill=signal!9, rounded corners=2.5pt, text=ink},
  gx/.style={draw=rule, fill=papertint, rounded corners=2.5pt, text=ink2}]
  \node[text=signal, anchor=west] at (0,1.0) {rank 0 $\cdot$ skalar};
  \node[bx, minimum width=0.7cm, minimum height=0.55cm] at (0.35,0.5) {\ttfamily 12};
  \node[text=ink3, anchor=west] at (0,-0.1) {shape ()};

  \node[text=signal, anchor=west] at (1.9,1.0) {rank 1 $\cdot$ vektor};
  \node[bx, minimum width=2.6cm, minimum height=0.55cm] at (3.2,0.5) {\ttfamily 12~~3~~6~~14~~7};
  \node[text=ink3, anchor=west] at (1.9,-0.1) {shape (5,)};

  \node[text=signal, anchor=west] at (5.5,1.0) {rank 2 $\cdot$ matriks};
  \node[bx, minimum width=2.5cm, minimum height=0.9cm, align=center] at (6.75,0.35)
    {\ttfamily 5~78~~2~34~0\\\ttfamily 6~79~~3~35~1};
  \node[text=ink3, anchor=west] at (5.5,-0.4) {shape (3, 5)};

  \node[text=signal, anchor=west] at (8.6,1.0) {rank 3};
  \node[gx, minimum width=1.7cm, minimum height=0.7cm] at (9.75,0.62) {};
  \node[gx, minimum width=1.7cm, minimum height=0.7cm] at (9.62,0.48) {};
  \node[bx, minimum width=1.7cm, minimum height=0.7cm] at (9.49,0.34) {\ttfamily matriks $\times$ n};
  \node[text=ink3, anchor=west] at (8.6,-0.4) {shape (3, 3, 5)};

  \draw[rule] (0,-0.8) -- (11.2,-0.8);
  \node[anchor=west, font=\bfseries\scriptsize, text=ink] at (0,-1.15) {MNIST sebagai tensor};
  \node[anchor=west, text=ink2] at (0,-1.5)
    {\ttfamily train\_images.ndim = 3 $\cdot$ shape (60000, 28, 28) $\cdot$ dtype uint8};
  \node[anchor=west, text=amber] at (0,-1.85)
    {sumbu ke-0 selalu sumbu sampel --- dan sumbu itulah yang dipotong jadi batch};
\end{tikzpicture}

}
\end{center}
\vskip 3pt{\color{ink3}\fontsize{7.4}{9.4}\selectfont Tensor menggeneralisasi matriks ke sebarang jumlah sumbu. TensorFlow diberi nama menurut benda ini.\par}
\begin{iband}[amber]
Jebakan istilah: \textbf{vektor 5 dimensi $\neq$ tensor 5 dimensi}. Yang pertama punya \hilite{satu sumbu berisi lima angka}; yang kedua punya \hilite{lima sumbu}. Salah baca di sini membuat pesan galat shape jadi tidak terbaca.
\end{iband}
\end{frame}
\note{Tanya ke peserta: shape sebuah batch 128 citra RGB 224x224 itu apa? Jawabannya (128, 224, 224, 3) --- dan itu latihan yang bagus.}

\framekicker{Bagian 2.2}
\begin{frame}[shrink=25]{Irisan tensor dan sumbu batch}
\icode{python}{listing 2.7-2.9 --- irisan}{listings/ch02-07.py}
\ioutput{listings/ch02-08.txt}
\begin{iband}[signal]
Sumbu ke-0 selalu \textbf{sumbu sampel}, dan karena model dilatih per potongan kecil, sumbu itu juga disebut \hilite{sumbu batch}. Setiap pesan galat shape yang akan Anda temui dimulai dari membaca sumbu ini.
\end{iband}
\end{frame}

\framekicker{Bagian 2.2.7-2.2.10}
\begin{frame}[shrink=25]{Data dunia nyata, dan bentuk tensornya}
\begin{itable}{>{\raggedright\arraybackslash}p{0.2068\textwidth}>{\raggedright\arraybackslash}p{0.0752\textwidth}>{\raggedright\arraybackslash}p{0.282\textwidth}>{\raggedright\arraybackslash}p{0.376\textwidth}}
\thead{Jenis data} & \thead{Rank} & \thead{Shape} & \thead{Contoh dari buku} \\ \midrule
\textbf{Vektor} & 2 & \inlinecode{(samples, features)} & 100.000 orang $\times$ (usia, jenis kelamin, penghasilan) $\rightarrow$ (100000, 3) \\
\textbf{Deret waktu} & 3 & \inlinecode{(samples, timesteps, features)} & 250 hari $\times$ 390 menit $\times$ 3 nilai $\rightarrow$ (250, 390, 3) \\
\textbf{Citra} & 4 & \inlinecode{(samples, h, w, channels)} & 128 citra RGB 256$\times$256 $\rightarrow$ (128, 256, 256, 3) \\
\textbf{Video} & 5 & \inlinecode{(samples, frames, h, w, channels)} & 4 klip $\times$ 240 bingkai $\times$ 144$\times$256 $\times$ RGB $\rightarrow$ (4, 240, 144, 256, 3) \\
\end{itable}
\begin{minipage}[t]{0.575\textwidth}
\begin{ibullets}
\item \textbf{Channels-last} \inlinecode{(\ldots{}, h, w, c)} -- kebiasaan TensorFlow dan JAX.
\item \textbf{Channels-first} \inlinecode{(\ldots{}, c, h, w)} -- kebiasaan PyTorch.
\item Keras 3 memakai \inlinecode{image\_data\_format} untuk memilih; salah setelan di sini \hilite{menghasilkan galat shape yang membingungkan}.
\end{ibullets}

\end{minipage}%
\hfill
\begin{minipage}[t]{0.395\textwidth}
\begin{minipage}[t]{1.0\textwidth}
\istat{106.168.320}{nilai dalam contoh video 60 detik itu}
\end{minipage}%


\vskip 4pt

\begin{minipage}[t]{1.0\textwidth}
\istat{425 MB}{ukurannya pada float32}
\end{minipage}%

\end{minipage}%
\end{frame}
\note{Data tabular dan data transaksi umumnya masuk baris pertama (vektor) atau kedua (deret waktu). Minta peserta menaksir shape datanya sendiri.}

\coursesection{03}{Gigi-giginya: operasi tensor}{Semuanya bermuara pada segenggam operasi -- dan semuanya punya arti geometris.}

\framekicker{Bagian 2.3}
\begin{frame}[shrink=25]{Satu lapis Dense = tiga operasi}
\icode{python}{inti sebuah lapis Dense}{listings/ch02-09.py}
\begin{minipage}[t]{0.485\textwidth}
\icode{python}{versi tervektorisasi}{listings/ch02-10.py}

\end{minipage}%
\hfill
\begin{minipage}[t]{0.485\textwidth}
\begin{minipage}[t]{0.494\textwidth}
\istat{0,02 s}{NumPy tervektorisasi, 1.000 iterasi}
\end{minipage}%
\hfill
\begin{minipage}[t]{0.494\textwidth}
\istat{2,45 s}{gelung Python naif, 1.000 iterasi}
\end{minipage}%

\end{minipage}%
\begin{iband}[signal]
Selisih \hilite{sekitar 100 kali} itu bukan soal bahasa. Versi NumPy melimpahkan pekerjaannya ke BLAS yang ditulis dalam C dan Fortran; di GPU, kode CUDA-nya tervektorisasi penuh.
\end{iband}
\end{frame}

\framekicker{Bagian 2.3.2}
\begin{frame}[shrink=25]{Broadcasting: bentuk kecil dipaskan ke bentuk besar}
\begin{center}
\adjustbox{max width=\linewidth,max totalheight=0.52\textheight}{%

\begin{tikzpicture}[font=\sffamily\tiny,
  ar/.style={-{Stealth[length=4pt]}, signal, line width=0.7pt}]
  \node[draw=signal!60, fill=signal!9, rounded corners=3pt, minimum width=1.7cm,
        minimum height=1.6cm, text=ink, align=center] (x) at (0,0) {X\\\ttfamily (32, 10)};
  \node[text=ink, font=\small] at (1.25,0) {$+$};
  \node[draw=rule, fill=papertint, rounded corners=3pt, minimum width=1.7cm,
        minimum height=0.42cm, text=ink2] (y) at (2.6,0) {\ttfamily y (10,)};
  \node[draw=rule, fill=papertint, rounded corners=3pt, minimum width=1.5cm,
        minimum height=0.42cm, text=ink2] (y2) at (5.0,0) {\ttfamily (1, 10)};
  \node[draw=lime!60, fill=limebr!12, rounded corners=3pt, minimum width=1.7cm,
        minimum height=1.6cm, text=ink, align=center] (Y) at (7.4,0) {Y\\\ttfamily (32, 10)};
  \draw[ar] (y) -- node[above, text=ink3, font=\tiny]{langkah 1} (y2);
  \draw[ar] (y2) -- node[above, text=ink3, font=\tiny]{langkah 2} (Y);
  \node[text=ink3, anchor=north] at (5.0,-0.35) {tambah sumbu};
  \node[text=ink3, anchor=north] at (7.4,-0.95) {diulang 32 kali};
  \node[text=amber, anchor=west] at (-1.0,-1.55)
    {tidak ada penggandaan memori sungguhan --- pengulangannya hanya algoritmis};
\end{tikzpicture}

}
\end{center}
\vskip 3pt{\color{ink3}\fontsize{7.4}{9.4}\selectfont Dua langkah: tambahkan sumbu sampai rank sama, lalu ulangi sepanjang sumbu baru itu.\par}
\icode{python}{listing 2.11 --- aturannya}{listings/ch02-11.py}
\ioutput{listings/ch02-12.txt}
\end{frame}
\note{Aturan umumnya: shape (a, b, ..., n, n+1, ..., m) berpasangan dengan (n, n+1, ..., m). Broadcasting adalah sumber bug diam yang paling sering --- bentuknya cocok, artinya tidak.}

\framekicker{Bagian 2.3.3-2.3.4}
\begin{frame}[shrink=25]{Hasil kali tensor dan reshape}
\begin{minipage}[t]{0.485\textwidth}
\icode{python}{matmul}{listings/ch02-13.py}

\end{minipage}%
\hfill
\begin{minipage}[t]{0.485\textwidth}
\icode{python}{reshape \& transpose}{listings/ch02-14.py}

\end{minipage}%
\begin{iband}[signal]
Reshape \hilite{tidak mengubah satu pun koefisien}; ia hanya menata ulang. Itulah yang terjadi pada \inlinecode{train\_images.reshape((60000, 784))} di contoh pertama tadi.
\end{iband}
\end{frame}

\framekicker{Bagian 2.3.5}
\begin{frame}[shrink=25]{Tafsir geometris -- dan mengapa aktivasi wajib ada}
\begin{itable}{>{\raggedright\arraybackslash}p{0.282\textwidth}>{\raggedright\arraybackslash}p{0.658\textwidth}}
\thead{Operasi} & \thead{Artinya secara geometris} \\ \midrule
Penjumlahan vektor & \textbf{Translasi} -- geser objek sejauh dan searah vektor itu. \\
Kali matriks rotasi & \textbf{Rotasi} sebesar sudut $\theta$. \\
Kali matriks diagonal & \textbf{Penskalaan} mendatar dan menegak. \\
Kali matriks sebarang & \textbf{Transformasi linear}. \\
Linear + translasi & \textbf{Transformasi afin} -- persis \inlinecode{y = W @ x + b}. \\
\end{itable}
\begin{iband}[rose]
Merangkai dua transformasi afin menghasilkan \textbf{satu} transformasi afin lagi: \inlinecode{affine2(affine1(x)) = (W2 @ W1) @ x + (W2 @ b1 + b2)}. Jadi tumpukan Dense tanpa aktivasi \hilite{diam-diam hanyalah satu model linear}, sedalam apa pun. Fungsi aktivasi seperti ReLU-lah yang membuat ruang hipotesisnya kaya.
\end{iband}
{\fontsize{9.4}{12.6}\selectfont\color{fg2}Gambaran yang dipakai Chollet: deep learning seperti \textbf{membuka remasan kertas}. Data yang terlipat rumit diluruskan sedikit demi sedikit oleh tiap lapis, sampai kelas-kelasnya bisa dipisah dengan bersih.\par}\vskip 4pt
\end{frame}
\note{Kalau ada satu hal dari bab 2 yang harus diingat manajer produk, ini: tanpa nonlinearitas, kedalaman tidak membeli apa pun.}

\coursesection{04}{Mesinnya: penurunan berbasis gradien}{Turunan, gradien, SGD, dan aturan rantai.}

\framekicker{Bagian 2.4}
\begin{frame}[shrink=25]{Lingkar pelatihan, dan langkah yang sulit}
\begin{isteps}
\item Ambil satu batch sampel \inlinecode{x} dan target \inlinecode{y\_true}.
\item Jalankan model pada \inlinecode{x} (\textbf{lintasan maju}) $\rightarrow$ \inlinecode{y\_pred}.
\item Hitung \textbf{rugi}: selisih antara \inlinecode{y\_pred} dan \inlinecode{y\_true}.
\item Perbarui bobot supaya ruginya turun sedikit.
\end{isteps}
\begin{iband}[amber]
Langkah 4 itulah yang sulit. Menyetel tiap koefisien satu per satu mustahil -- jaringan modern punya \hilite{jutaan sampai miliaran parameter}. Penurunan gradien menyelesaikannya sekaligus.
\end{iband}
\begin{minipage}[t]{0.485\textwidth}
{\fontsize{9.4}{12.6}\selectfont\color{fg2}\textbf{Turunan} adalah kemiringan hampiran linear setempat: untuk $\varepsilon$\_x cukup kecil, \inlinecode{f(x + $\varepsilon$\_x) $\approx$ y + a$\cdot$$\varepsilon$\_x}.\par}\vskip 4pt
\begin{ibullets}
\item \inlinecode{a} negatif $\rightarrow$ menaikkan \inlinecode{x} \textbf{menurunkan} \inlinecode{f(x)}.
\item \inlinecode{a} positif $\rightarrow$ menaikkan \inlinecode{x} \textbf{menaikkan} \inlinecode{f(x)}.
\item Untuk mengecilkan \inlinecode{f}, geser \inlinecode{x} \hilite{berlawanan arah turunannya}.
\end{ibullets}

\end{minipage}%
\hfill
\begin{minipage}[t]{0.485\textwidth}
{\fontsize{9.4}{12.6}\selectfont\color{fg2}\textbf{Gradien} adalah turunan untuk operasi tensor. Ia tensor sebentuk \inlinecode{W}, dan tiap koefisiennya menyatakan arah dan besar perubahan rugi bila koefisien \inlinecode{W} itu digeser.\par}\vskip 4pt
\begin{iband}[signal]
\inlinecode{W1 = W0 - step * grad(f(W0), W0)}
\end{iband}

\end{minipage}%
\end{frame}

\framekicker{Bagian 2.4.3}
\begin{frame}[shrink=25]{SGD mini-batch, dan mengapa tidak diselesaikan secara analitis}
{\fontsize{9.4}{12.6}\selectfont\color{fg2}Minimum ada di tempat turunannya nol. Tetapi menyelesaikan \inlinecode{grad(f(W), W) = 0} secara analitis \hilite{tidak terjangkau} untuk jaringan berparameter jutaan. Maka: iterasi.\par}\vskip 4pt
\begin{isteps}
\item Ambil batch acak \inlinecode{x}, \inlinecode{y\_true}. (Kata \textit{stochastic} datang dari \textbf{acak} ini.)
\item Lintasan maju $\rightarrow$ \inlinecode{y\_pred}.
\item Hitung rugi.
\item \textbf{Lintasan mundur} $\rightarrow$ gradien rugi terhadap parameter.
\item \inlinecode{W -= learning\_rate * gradient}.
\end{isteps}
\begin{minipage}[t]{0.3253\textwidth}
\begin{icardwarn}
\cardh{Learning rate terlalu kecil}
Kekonvergenan lambat; banyak langkah untuk sedikit kemajuan.
\end{icardwarn}
\end{minipage}%
\hfill
\begin{minipage}[t]{0.3253\textwidth}
\begin{icardbad}
\cardh{Terlalu besar}
Pembaruan jadi \hilite{acak}; rugi melompat-lompat dan tidak turun.
\end{icardbad}
\end{minipage}%
\hfill
\begin{minipage}[t]{0.3253\textwidth}
\begin{icardgood}
\cardh{Momentum}
Perbarui berdasar gradien \textbf{sekarang dan pembaruan sebelumnya} -- seperti bola menggelinding, cukup laju untuk melewati cekungan dangkal.
\end{icardgood}
\end{minipage}%
{\fontsize{9.4}{12.6}\selectfont\color{fg2}Varian yang memperbaiki kekonvergenan disebut \textbf{optimizer}: SGD dengan momentum, Adagrad, RMSprop, Adam.\par}\vskip 4pt
\end{frame}

\framekicker{Bagian 2.4.4-2.4.5}
\begin{frame}[shrink=25]{Aturan rantai di atas graf komputasi = backpropagation}
\begin{center}
\adjustbox{max width=\linewidth,max totalheight=0.52\textheight}{%

\begin{tikzpicture}[font=\sffamily\tiny,
  bx/.style={draw=rule, fill=papertint, rounded corners=3pt, minimum width=1.5cm,
             minimum height=0.5cm, text=ink2},
  ax/.style={draw=signal!60, fill=signal!9, rounded corners=3pt, minimum width=1.7cm,
             minimum height=0.5cm, text=ink},
  fw/.style={-{Stealth[length=4pt]}, signal, line width=0.7pt},
  bw/.style={-{Stealth[length=4pt]}, amberbr, line width=0.8pt}]

  \node[text=signal, anchor=west] at (0,0.7) {lintasan maju --- hitung nilai};
  \node[bx] (x)  at (0.75,0)  {\ttfamily x};
  \node[ax] (x1) at (2.9,0)   {\ttfamily x1 = W$\cdot$x};
  \node[ax] (x2) at (5.1,0)   {\ttfamily x2 = x1+b};
  \node[ax] (y)  at (7.4,0)   {\ttfamily y = relu(x2)};
  \node[draw=rose!70, fill=rosebr!14, rounded corners=3pt, minimum width=1.4cm,
        minimum height=0.5cm, text=ink] (l) at (9.6,0) {\ttfamily loss};
  \draw[fw] (x) -- (x1); \draw[fw] (x1) -- (x2); \draw[fw] (x2) -- (y); \draw[fw] (y) -- (l);

  \draw[rule] (0,-0.55) -- (10.5,-0.55);
  \node[text=amber, anchor=west] at (0,-0.9)
    {lintasan mundur --- balik arah sisi, kalikan turunan sepanjang lintasan};
  \draw[bw] (9.0,-1.35) -- (8.1,-1.35);
  \draw[bw] (6.8,-1.35) -- (5.8,-1.35);
  \draw[bw] (4.5,-1.35) -- (3.6,-1.35);
  \draw[bw] (2.3,-1.35) -- (1.4,-1.35);
  \node[text=amber, font=\ttfamily\tiny] at (8.55,-1.1)  {$\partial$loss/$\partial$y};
  \node[text=amber, font=\ttfamily\tiny] at (6.3,-1.1)   {$\partial$y/$\partial$x2};
  \node[text=amber, font=\ttfamily\tiny] at (4.05,-1.1)  {$\partial$x2/$\partial$x1};
  \node[text=amber, font=\ttfamily\tiny] at (1.85,-1.1)  {$\partial$x1/$\partial$W};

  \node[draw=amber!45, fill=amberbr!8, rounded corners=4pt, minimum width=10.2cm,
        minimum height=0.95cm, anchor=north west, align=left] at (0,-1.75) {};
  \node[anchor=west, text=ink, font=\ttfamily\tiny] at (0.25,-2.05)
    {grad(loss, W) = $\partial$loss/$\partial$y $\times$ $\partial$y/$\partial$x2 $\times$ $\partial$x2/$\partial$x1 $\times$ $\partial$x1/$\partial$W};
  \node[anchor=west, text=ink3] at (0.25,-2.4)
    {aturan rantai, dijalankan mundur di atas graf --- itulah backpropagation};
\end{tikzpicture}

}
\end{center}
\vskip 3pt{\color{ink3}\fontsize{7.4}{9.4}\selectfont Graf komputasi adalah graf berarah tanpa siklus. Ia membuat perhitungan bisa diperlakukan sebagai data -- struktur yang terbaca mesin.\par}
\begin{ibullets}
\item Bila ada \textbf{beberapa lintasan} dari simpul \inlinecode{a} ke \inlinecode{b}, sumbangan tiap lintasan \hilite{dijumlahkan}.
\item Kerangka kerja modern menjalankan ini sendiri: itulah \textbf{automatic differentiation}. Anda tidak pernah menulis backprop dengan tangan.
\end{ibullets}
\end{frame}
\note{Kalau peserta pernah menurunkan backprop manual di kuliah S1, katakan bahwa yang mereka kerjakan dulu itu persis ini --- hanya saja kini dikerjakan oleh graf.}

\framekicker{Bagian 2.4.5}
\begin{frame}[shrink=25]{GradientTape: gradien dalam tiga baris}
\icode{python}{listing 2.19-2.20 --- autodiff}{listings/ch02-15.py}
\ioutput{listings/ch02-16.txt}
\begin{iband}[signal]
Contoh kedua bukan sekadar pamer: \inlinecode{position = 4.9$\cdot$t$^{2}$} adalah jatuh bebas, dan turunan keduanya memang \hilite{percepatan gravitasi}. Autodiff mengembalikan 9,8 tanpa pernah diberi tahu rumusnya.
\end{iband}
\end{frame}

\coursesection{05}{Menulis ulang dari nol}{Tanpa fit(), tanpa Dense bawaan. Hasilnya harus sama.}

\framekicker{Bagian 2.6.1}
\begin{frame}[shrink=25]{NaiveDense dan NaiveSequential}
\icode{python}{listing 2.21-2.22 --- lapis dan model}{listings/ch02-17.py}
\end{frame}

\framekicker{Bagian 2.6.2-2.6.4}
\begin{frame}[shrink=25]{Satu langkah pelatihan, dan gelungnya}
\icode{python}{listing 2.24-2.26 --- lingkar pelatihan}{listings/ch02-18.py}
\begin{iband}[signal]
Empat langkah pada slide lingkar pelatihan tadi kini terlihat sebagai empat baris: \hilite{maju, rugi, gradien, perbarui}.
\end{iband}
\end{frame}

\framekicker{Bagian 2.6.5}
\begin{frame}[shrink=25]{Buktinya: hasilnya memang sama}
\icode{python}{listing 2.27 --- menilai model}{listings/ch02-19.py}
\ioutput{listings/ch02-20.txt}
\begin{iband}[amber]
Sengaja \hilite{lebih rendah dari 97,8\%}. Bedanya bukan sihir Keras: versi ini memakai SGD polos dengan learning rate 1e-3, sedangkan yang pertama memakai \textbf{Adam}. Itu justru pelajarannya -- pilihan optimizer berdampak sebesar itu.
\end{iband}
\end{frame}
\note{Ini slide yang paling sering salah dibaca. Tegaskan: turunnya akurasi bukan bukti bahwa implementasi manualnya salah.}

\framekicker{Ringkasan}
\begin{frame}[shrink=25]{Yang wajib terbawa dari bab 2}
\begin{isteps}
\item \textbf{Tensor} = wadah angka dengan \textbf{rank, shape, dtype}. Sumbu ke-0 adalah sumbu sampel alias sumbu batch.
\item \textbf{Operasi tensor} punya arti geometris. Dense tanpa aktivasi hanya transformasi afin -- dan tumpukannya tetap afin.
\item \textbf{Belajar} = mencari nilai bobot yang mengecilkan rugi pada data latih.
\item \textbf{SGD mini-batch} memperbarui bobot dari gradien pada batch acak.
\item \textbf{Aturan rantai + autodiff} = backpropagation; graf komputasi yang mengerjakannya.
\item \textbf{Loss} menakar keberhasilan tugas; \textbf{optimizer} menentukan varian penurunan gradiennya.
\end{isteps}
\vskip 4pt
\reslink{NOTEBOOK}{01\_mnist\_pertama.ipynb}{../../course-slides/notebooks/ch02/01\_mnist\_pertama.ipynb}
\reslink{NOTEBOOK}{04\_mnist\_dari\_nol.ipynb}{../../course-slides/notebooks/ch02/04\_mnist\_dari\_nol.ipynb}
\reslink{BAB BERIKUT}{Bab 3 --- Kerangka kerja}{../ch03/index.html}
\vskip 2pt
\end{frame}

\end{document}
